\documentclass[multi=tikzpicture,border=9pt]{standalone}
\usepackage[UTF8]{ctex}
\usepackage{amsmath,amssymb}
\usepackage{tikz}
\usetikzlibrary{arrows.meta,calc,patterns,decorations.pathmorphing}
\definecolor{main}{HTML}{17365D}
\definecolor{accent}{HTML}{AC3931}
\tikzset{every picture/.style={font=\small,line width=.75pt},
>={Stealth[length=2mm]},arr/.style={->},box/.style={draw=main,rounded corners=2pt,
fill=main!4,text width=2.7cm,minimum height=.9cm,align=center},
note/.style={align=center,text width=8cm}}
\begin{document}
\begin{tikzpicture}
\draw[thick] (0,0) rectangle (3,3);\draw (0,3)--(.8,3.7)--(3.8,3.7)--(3,3);\draw (3,0)--(3.8,.7)--(3.8,3.7);\draw[arr] (-.6,0)--(4.5,0) node[right]{$x$};\draw[arr] (0,-.3)--(0,4.2) node[above]{$y$};\fill (0.6,0.7) circle(2pt);\draw[arr,main] (0.6,0.7)--(1.25,0.7);\fill (2,1.5) circle(2pt);\draw[arr,main] (2,1.5)--(1.35,1.5);\fill (0.8,2.2) circle(2pt);\draw[arr,main] (0.8,2.2)--(1.4500000000000002,2.2);\node[below] at (1.5,0){$l$};\node[right] at (3.9,2){$l$};\node[note] at (1.7,-1.25){箱体静止系仅有 $\pm v\,\mathbf e_x$；左右壁发生碰撞，其他原壁面没有碰撞。};
\end{tikzpicture}
\begin{tikzpicture}
\node[box] (b0) at (0,0) {单次壁面冲量\\$2p_*$};
\node[box] (b1) at (5,0) {同壁周期\\$\tau_0=2l/v$};
\node[box] (b2) at (5,-2) {面积\\$A_0=l^2$};
\node[box] (b3) at (0,-2) {$p_0=n\gamma_vm_0v^2$};
\draw[arr] (b0)--(b1);\draw[arr] (b1)--(b2);\draw[arr] (b2)--(b3);\node[note] at (2.5,-3.35) {沿y观察：力与面积同除 $\gamma_u$；沿x观察：冲量与周期同乘 $\gamma_u$。};
\end{tikzpicture}
\begin{tikzpicture}
\draw[thick] (0,0)--(0,2.5);\draw[thick] (4,0)--(4,2.5);\draw[arr,accent] (1,2.7)--(0,2.7);\draw[arr,accent] (5,2.7)--(4,2.7);\node[above] at (2.5,2.7){两壁速度 $-u$};\draw[arr,main] (1,.8)--(2.7,.8) node[midway,below]{$v_+^{\prime}$};\draw[arr,main] (3,1.6)--(1.3,1.6) node[midway,above]{$v_-^{\prime}$};\draw[<->] (0,-.5)--(4,-.5) node[midway,below]{$l/\gamma_u$};\node[note] at (2,-1.7){追及时间按粒子相对壁面速度求；两段之和 $\tau^{\prime}=\gamma_u\,2l/v$。箭头方向为 $u<v$ 的示意。};
\end{tikzpicture}
\begin{tikzpicture}
\draw[arr] (0,0)--(5,0) node[right]{$\mathbf e$};\draw[arr] (0,0)--(0,3.2) node[above]{$\mathbf t$};\draw[main,very thick,arr] (0,0)--(3,-1.8);\node[anchor=west,main] at (3.2,-1.8){$\mathbf A^{\prime}/A_0$};\draw[accent,arr] (1.5,2)--(4.5,.2);\node[anchor=west,accent] at (4.65,.2){$\mathbf F^{\prime}/F_0$};\draw[dashed] (3,0)--(3,-1.8)--(0,-1.8);\node[above] at (1.5,0){$\cos\theta$};\node[left] at (0,-.9){$-\sin\theta/\gamma_u$};\node[note] at (2,-3){两个矢量方向相同，长度缩放均为 $K=\sqrt{\cos^2\theta+\sin^2\theta/\gamma_u^2}$。};
\end{tikzpicture}
\begin{tikzpicture}
\node[box] (b0) at (0,0) {指定同一移动壁面};
\node[box] (b1) at (5,0) {变换冲量与时间};
\node[box] (b2) at (5,-2) {变换有向面积};
\node[box] (b3) at (0,-2) {法向力/面积\\仍为 $p_0$};
\draw[arr] (b0)--(b1);\draw[arr] (b1)--(b2);\draw[arr] (b2)--(b3);\node[note] at (2.5,-3.35) {固定空间平面的动量通量可含整体输运项；原来不碰撞的侧壁不会新增碰撞。};
\end{tikzpicture}
\end{document}
